\documentclass[14pt,a4paper]{extarticle}

% ─── Кодировка и язык ────────────────────────────────────────────────────────
\usepackage[utf8]{inputenc}
\usepackage[T2A]{fontenc}
\usepackage[russian]{babel}

% ─── Математические символы (для \checkmark) ─────────────────────────────────
\usepackage{amssymb}

% ─── Поля страницы ───────────────────────────────────────────────────────────
\usepackage[left=30mm,right=15mm,top=20mm,bottom=20mm]{geometry}

% ─── Межстрочный интервал ────────────────────────────────────────────────────
\usepackage{setspace}
\onehalfspacing

% ─── Отступ первого абзаца ───────────────────────────────────────────────────
\usepackage{indentfirst}
\setlength{\parindent}{1.25cm}

% ─── Таблицы ─────────────────────────────────────────────────────────────────
\usepackage{booktabs}
\usepackage{array}
\usepackage{longtable}
\usepackage{tabularx}

% ─── Листинги кода ───────────────────────────────────────────────────────────
\usepackage{listings}
\usepackage{xcolor}

\definecolor{codebg}{HTML}{F5F5F5}
\definecolor{codeframe}{HTML}{CCCCCC}
\definecolor{codecomment}{HTML}{5C8A5C}
\definecolor{codestring}{HTML}{A31515}
\definecolor{codekeyword}{HTML}{0000CC}

% Определение синтаксиса YAML для пакета listings
\lstdefinelanguage{yaml}{
	keywords={services, image, container_name, restart, privileged, environment, volumes, networks, ipv4_address, depends_on, condition, driver, driver_opts, type, o, device, healthcheck, test, interval, timeout, retries, start_period},
	keywordstyle=\color{codekeyword}\bfseries,
	sensitive=false,
	comment=[l]{\#},
	commentstyle=\color{codecomment},
	stringstyle=\color{codestring},
	morestring=[b]',
	morestring=[b]"
}

\lstset{
	backgroundcolor=\color{codebg},
	basicstyle=\ttfamily\small,
	breakatwhitespace=false,
	breaklines=true,
	captionpos=b,
	commentstyle=\color{codecomment},
	frame=single,
	rulecolor=\color{codeframe},
	keepspaces=true,
	keywordstyle=\color{codekeyword}\bfseries,
	numbers=left,
	numbersep=8pt,
	numberstyle=\tiny\color{gray},
	showspaces=false,
	showstringspaces=false,
	showtabs=false,
	stringstyle=\color{codestring},
	tabsize=2,
	extendedchars=true,
	inputencoding=utf8,
	literate=
	{а}{{\cyrа}}1 {б}{{\cyrb}}1 {в}{{\cyrv}}1 {г}{{\cyrg}}1 {д}{{\cyrd}}1
	{е}{{\cyre}}1 {ё}{{\cyryo}}1 {ж}{{\cyrzh}}1 {з}{{\cyrz}}1 {и}{{\cyri}}1
	{й}{{\cyrishrt}}1 {к}{{\cyrk}}1 {л}{{\cyrl}}1 {м}{{\cyrm}}1 {н}{{\cyrn}}1
	{о}{{\cyro}}1 {п}{{\cyrp}}1 {р}{{\cyrr}}1 {с}{{\cyrs}}1 {т}{{\cyrt}}1
	{у}{{\cyru}}1 {ф}{{\cyrf}}1 {х}{{\cyrh}}1 {ц}{{\cyrc}}1 {ч}{{\cyrch}}1
	{ш}{{\cyrsh}}1 {щ}{{\cyrshch}}1 {ъ}{{\cyrhrdsn}}1 {ы}{{\cyrery}}1
	{ь}{{\cyrsftsn}}1 {э}{{\cyrerev}}1 {ю}{{\cyryu}}1 {я}{{\cyrya}}1
	{А}{{\CYRA}}1 {Б}{{\CYRB}}1 {В}{{\CYRV}}1 {Г}{{\CYRG}}1 {Д}{{\CYRD}}1
	{Е}{{\CYRE}}1 {Ё}{{\CYRYO}}1 {Ж}{{\CYRZH}}1 {З}{{\CYRZ}}1 {И}{{\CYRI}}1
	{Й}{{\CYRISHRT}}1 {К}{{\CYRK}}1 {Л}{{\CYRL}}1 {М}{{\CYRM}}1 {Н}{{\CYRN}}1
	{О}{{\CYRO}}1 {П}{{\CYRP}}1 {Р}{{\CYRR}}1 {С}{{\CYRS}}1 {Т}{{\CYRT}}1
	{У}{{\CYRU}}1 {Ф}{{\CYRF}}1 {Х}{{\CYRH}}1 {Ц}{{\CYRC}}1 {Ч}{{\CYRCH}}1
	{Ш}{{\CYRSH}}1 {Щ}{{\CYRSHCH}}1 {Ъ}{{\CYRHRDSN}}1 {Ы}{{\CYRERY}}1
	{Ь}{{\CYRSFTSN}}1 {Э}{{\CYREREV}}1 {Ю}{{\CYRYU}}1 {Я}{{\CYRYA}}1
}

% ─── Гиперссылки ─────────────────────────────────────────────────────────────
\usepackage[hidelinks]{hyperref}

% ─── Колонтитулы ─────────────────────────────────────────────────────────────
\usepackage{fancyhdr}
\pagestyle{fancy}
\fancyhf{}
\fancyhead[R]{\small Курсовой проект | Администрирование ОС Linux | ЮФУ 2026}
\fancyfoot[C]{\thepage}
\renewcommand{\headrulewidth}{0.4pt}

% ─── Заголовки секций ────────────────────────────────────────────────────────
\usepackage{titlesec}
\titleformat{\section}{\large\bfseries}{}{0em}{\thesection\quad}
\titleformat{\subsection}{\normalsize\bfseries}{}{0em}{\thesubsection\quad}

% ─── Переопределение счётчика страниц для титульника ─────────────────────────
\usepackage{etoolbox}

\begin{document}
	
	% ════════════════════════════════════════════════════════════════════════════════
	%  ТИТУЛЬНАЯ СТРАНИЦА
	% ════════════════════════════════════════════════════════════════════════════════
	\thispagestyle{empty}
	\begin{center}
		{\normalsize
			Министерство науки и высшего образования Российской Федерации\\[4pt]
			Южный федеральный университет\\[4pt]
			Физический факультет\\[40pt]
		}
		
		{\Large\bfseries КУРСОВОЙ ПРОЕКТ}\\[10pt]
		{\normalsize по дисциплине «Администрирование ОС Linux»}\\[30pt]
		
		{\large\bfseries Тема: Использование NFS-сервера для сетевого хранения данных}\\[50pt]
		
		\begin{flushleft}
			\hspace{6cm}
			\begin{minipage}{0.5\textwidth}
				\textbf{Авторы (Группа):}\\[4pt]
				Григорян А.\,А.\\
				Карлашов А.\,В.\\
				Кроливец И.\,С.\\[20pt]
				\textbf{Проверил:}\\[4pt]
				Тимошенко Павел Евгеньевич
			\end{minipage}
		\end{flushleft}
		
		\vfill
		{\normalsize Ростов-на-Дону, 2026}
	\end{center}
	
	\newpage
	
	% ════════════════════════════════════════════════════════════════════════════════
	%  СОДЕРЖАНИЕ
	% ════════════════════════════════════════════════════════════════════════════════
	\thispagestyle{fancy}
	\tableofcontents
	
	\newpage
	
	% ════════════════════════════════════════════════════════════════════════════════
	%  ВВЕДЕНИЕ
	% ════════════════════════════════════════════════════════════════════════════════
	\section*{Введение}
	\addcontentsline{toc}{section}{Введение}
	
	Современный этап развития корпоративной ИТ-инфраструктуры предъявляет всё более
	высокие требования к централизованному хранению и совместному использованию
	файловых ресурсов. В средах с большим числом узлов --- будь то физические серверы,
	виртуальные машины или контейнеми --- каждый узел не должен хранить собственную копию
	данных; необходима единая точка хранения, доступная всем участникам инфраструктуры
	одновременно.
	
	Протокол NFS (Network File System), разработанный компанией Sun Microsystems в
	1984 году и прошедший эволюцию вплоть до версии NFSv4.2, остаётся отраслевым
	стандартом сетевых файловых систем для Unix-подобных сред. Его нативная интеграция с
	ядром Linux, высокая производительность при работе с большим числом файлов и богатые
	возможности разграничения прав доступа делают NFS оптимальным выбором для задач
	горизонтального масштабирования хранилищ и организации совместной работы контейнеров
	с общими данными.
	
	Однако ручная настройка NFS-сервера неизбежно порождает «дрейф конфигураций»
	(configuration drift) и ошибки воспроизводимости. Применение методологии
	Infrastructure as Code (IaC) совместно с инструментами контейнеризации Docker
	позволяет полностью автоматизировать жизненный цикл развёртывания, устранив
	человеческий фактор.
	
	\textbf{Цель данного проекта} --- комплексное проектирование, развёртывание, настройка
	и защита NFS-сервера в контейнеризированной среде под управлением ОС Linux, а также
	обеспечение безопасного сетевого доступа к нему для множества клиентских контейнеров.
	
	\newpage
	
	% ════════════════════════════════════════════════════════════════════════════════
	%  1. РАСПРЕДЕЛЕНИЕ РОЛЕЙ
	% ════════════════════════════════════════════════════════════════════════════════
	\section{Распределение ролей в группе и индивидуальный вклад}
	
	Разработка, тестирование и обеспечение безопасности сетевого файлового хранилища
	выполнялись проектной группой из трёх человек с чётким разграничением обязанностей и
	зон инженерной ответственности.
	
	\begin{table}[h!]
		\centering
		\small
		\begin{tabularx}{\textwidth}{|l|l|X|}
			\hline
			\textbf{Участник} & \textbf{Роль} & \textbf{Основные задачи} \\
			\hline
			Григорян Ашот & Инженер по автоматизации & Docker Compose, bootstrap.sh, Dockerfile, Makefile \\
			\hline
			Карлашов Андрей & Системный администратор / SRE & Настройка NFS, монтирование, healthcheck, HA-тесты \\
			\hline
			Кроливец Иван & Инженер по безопасности & Hardening, UFW, Prometheus, Grafana, CI/CD-пайплайн \\
			\hline
		\end{tabularx}
		\caption{Распределение ролей в проектной группе}
	\end{table}
	
	\subsection{Инженер по автоматизации инфраструктуры --- Григорян Ашот}
	
	В зону ответственности инженера по автоматизации входила разработка декларативного
	описания всей инфраструктуры и автоматизация её жизненного цикла:
	
	\begin{itemize}
		\item Проектирование структуры Git-репозитория и настройка ветвления
		(\texttt{main}, \texttt{feature/*}, \texttt{gh-pages});
		\item Разработка главного оркестрационного манифеста \texttt{docker-compose.yml}
		с описанием NFS-сервера, клиентских контейнеров и изолированных виртуальных сетей;
		\item Написание скрипта \texttt{bootstrap.sh} первичной подготовки ОС: установка
		зависимостей, создание директорий экспорта, настройка прав доступа;
		\item Сборка оптимизированных Dockerfile с применением принципа наименьших
		привилегий (non-root user) и многоэтапной сборки (multi-stage builds);
		\item Автоматизация управления стендом через Makefile: цели
		\texttt{up}, \texttt{down}, \texttt{test}, \texttt{clean}.
	\end{itemize}
	
	\subsection{Системный администратор / SRE --- Карлашов Андрей}
	
	Системный администратор сфокусировался на конфигурации сервисного уровня и
	обеспечении отказоустойчивости:
	
	\begin{itemize}
		\item Настройка NFS-сервера (\texttt{nfs-kernel-server} / NFS Ganesha): конфигурация
		файла \texttt{/etc/exports}, опций \texttt{rw}, \texttt{sync},
		\texttt{root\_squash}, \texttt{fsid};
		\item Конфигурация нескольких клиентских контейнеров с автоматическим монтированием
		NFS-шары через Docker Volume (\texttt{driver: local, type: nfs});
		\item Написание и интеграция скриптов healthcheck для проверки доступности
		NFS-экспорта;
		\item Проведение HA-тестов: имитация падения сервера, замер времени восстановления
		доступа к данным клиентами;
		\item Проверка блокировок файлов при одновременной записи из нескольких контейнеров.
	\end{itemize}
	
	\subsection{Инженер по мониторингу и безопасности --- Кроливец Иван}
	
	Инженер по безопасности обеспечивал многоуровневую защиту хранилища и наблюдаемость
	системы:
	
	\begin{itemize}
		\item Разработка скрипта \texttt{hardening.sh}: конфигурация UFW по принципу
		«запрещено всё, кроме явно разрешённого»; открыты только порты SSH (22),
		NFS (2049) для доверенных подсетей;
		\item Изоляция сети: NFS-трафик строго ограничен виртуальной сетью
		\texttt{storage-net}, внешний доступ к порту 2049 заблокирован;
		\item Развёртывание стека мониторинга Prometheus + Grafana с экспортером
		\texttt{node\_exporter} и cAdvisor для сбора метрик контейнеров;
		\item Разработка дашбордов Grafana: нагрузка на NFS-шару, I/O дисковой подсистемы,
		состояние контейнеров;
		\item Настройка конвейера CI/CD в GitHub Actions: линтеры ShellCheck, yamllint,
		Hadolint; блокировка слияния при провале тестов.
	\end{itemize}
	
	\newpage
	
	% ════════════════════════════════════════════════════════════════════════════════
	%  2. ТЕОРЕТИЧЕСКИЕ ОСНОВЫ
	% ════════════════════════════════════════════════════════════════════════════════
	\section{Теоретические основы технологии NFS}
	
	NFS (Network File System) --- клиент-серверный протокол, позволяющий монтировать
	удалённые файловые системы как локальные. В ядре Linux реализована поддержка NFSv3 и
	NFSv4/4.1/4.2. Архитектура NFS строится на модели RPC (Remote Procedure Call) с
	транспортом по протоколам TCP и UDP.
	
	Версия NFSv4 принципиально отличается от предшественников: она использует единственный
	порт 2049/TCP, обеспечивает интегрированную безопасность через Kerberos или
	RPCSEC\_GSS, поддерживает сессионную модель с управлением состоянием соединения и
	атомарные операции аренды (lease). NFSv4.1 добавляет параллельное расположение данных
	(pNFS), а NFSv4.2 --- серверную копию файлов без участия клиента (server-side copy).
	
	\textbf{Ключевые параметры экспорта NFS}, задаваемые в файле \texttt{/etc/exports}:
	
	\begin{itemize}
		\item \texttt{rw / ro} --- разрешение записи или только чтения для клиента;
		\item \texttt{sync} --- гарантия записи данных на диск до подтверждения клиенту
		(важно для целостности данных);
		\item \texttt{async} --- более высокая производительность, но риск потери данных при
		сбое;
		\item \texttt{root\_squash} --- перемаппинг root-пользователя клиента в анонимный
		UID/GID (безопасность);
		\item \texttt{no\_root\_squash} --- отключение маппинга (используется только в
		доверенных средах);
		\item \texttt{fsid=0} --- обязателен для NFSv4, определяет корень псевдофайловой
		системы экспорта.
	\end{itemize}
	
	\textbf{Сравнение NFS с альтернативными решениями.}
	При выборе технологии хранения была проведена сравнительная оценка трёх подходов:
	
	\begin{table}[h!]
		\centering
		\small
		\begin{tabularx}{\textwidth}{|l|X|X|l|}
			\hline
			\textbf{Технология} & \textbf{Преимущества} & \textbf{Недостатки} & \textbf{Вывод} \\
			\hline
			NFS &
			Нативная поддержка в Linux, скорость, простая интеграция с Docker &
			Ограниченная отказоустойчивость, слабая защита без Kerberos &
			Оптимален \\
			\hline
			Samba (SMB) &
			Кросс-платформенность, интеграция с AD/LDAP &
			Выше накладные расходы в Linux, сложнее интеграция с Docker &
			Избыточен \\
			\hline
			GlusterFS &
			Распределённая архитектура, репликация, масштабирование &
			Высокая сложность настройки, требует от 3 узлов кластера &
			Избыточен \\
			\hline
		\end{tabularx}
		\caption{Сравнение технологий сетевого файлового хранилища}
	\end{table}
	
	\newpage
	
	% ════════════════════════════════════════════════════════════════════════════════
	%  3. ПРОЕКТИРОВАНИЕ АРХИТЕКТУРЫ
	% ════════════════════════════════════════════════════════════════════════════════
	\section{Проектирование архитектуры и сетевая изоляция}
	
	Архитектура стенда построена по принципу минимальной связности компонентов. Все
	сервисы развёрнуты в Docker-контейнерах и взаимодействуют исключительно через
	изолированные виртуальные сети.
	
	\textbf{Состав кластера:}
	
	\begin{description}
		\item[\texttt{nfs-server}] Основной сервис хранения данных. Экспортирует директорию
		\texttt{/exports/shared} по протоколу NFSv4 на порту 2049. Примонтированный
		именованный том \texttt{nfs\_data} обеспечивает персистентность данных.
		
		\item[\texttt{nfs-client-1}] Первый клиент. Автоматически монтирует NFS-шару через
		Docker Volume (\texttt{driver: local, type: nfs4}) при старте контейнера.
		Имитирует сервис приложения.
		
		\item[\texttt{nfs-client-2}] Второй клиент. Монтирует ту же NFS-шару, демонстрируя
		совместный доступ к данным из нескольких контейнеров одновременно.
		
		\item[\texttt{monitoring}] Контейнер Prometheus + Grafana для сбора метрик нагрузки
		на NFS-сервер и дисковой подсистемы.
	\end{description}
	
	\textbf{Сетевая topology.}
	Безопасность сетевого контура реализована через виртуализацию Docker.
	Созданы две изолированные сети:
	
	\begin{itemize}
		\item \texttt{storage-net} (bridge) --- приватная сеть для NFS-трафика. NFS-сервер и
		оба клиента подключены исключительно к этой сети. Внешний доступ к порту 2049
		отсутствует.
		\item \texttt{monitoring-net} (bridge) --- отдельная сеть для Prometheus/Grafana.
		Дашборды Grafana доступны на порту 3000 хост-машины с базовой
		HTTP-аутентификацией.
	\end{itemize}
	
	Сетевое разграничение гарантирует, что NFS-трафик (в том числе данные файлов) не
	может быть перехвачен другими контейнерами и не покидает виртуальный сегмент
	\texttt{storage-net}.
	
	\newpage
	
	% ════════════════════════════════════════════════════════════════════════════════
	%  4. АВТОМАТИЗАЦИЯ
	% ════════════════════════════════════════════════════════════════════════════════
	\section{Автоматизация развёртывания и концепция IaC}
	
	Развёртывание инфраструктуры реализовано в рамках парадигмы IaC с помощью
	декларативного инструмента Docker Compose. Это полностью исключает дрейф конфигураций
	и гарантирует идентичность сред разработки, тестирования и промышленной эксплуатации.
	
	\subsection{Декларативное описание стека}
	
	Ниже приведён фрагмент конфигурационного файла \texttt{docker-compose.yml},
	описывающий NFS-сервер и одного из клиентов:
	
	\begin{lstlisting}[language=yaml, caption={Фрагмент docker-compose.yml}]
		services:
		nfs-server:
		image: erichough/nfs-server:latest
		container_name: nfs-server
		restart: always
		privileged: true
		environment:
		NFS_EXPORT_0: '/exports/shared *(rw,sync,no_subtree_check,fsid=0)'
		volumes:
		- nfs_data:/exports/shared
		networks:
		storage-net:
		ipv4_address: 172.20.0.10
		
		nfs-client-1:
		image: alpine:3.19
		container_name: nfs-client-1
		restart: always
		depends_on:
		nfs-server:
		condition: service_healthy
		volumes:
		- type: volume
		source: nfs_shared
		target: /mnt/shared
		networks:
		- storage-net
	\end{lstlisting}
	
	\subsection{Автоматизация жизненного цикла окружения}
	
	Ниже приведена реализация скрипта первичной инициализации хост-системы:
	
	\begin{lstlisting}[language=bash, caption={Скрипт bootstrap.sh}]
		#!/bin/bash
		# Host initialization script
		set -euo pipefail
		
		EXPORT_DIR="${EXPORT_DIR:-/opt/nfs/shared}"
		BACKUP_DIR="${BACKUP_DIR:-/opt/nfs/backups}"
		
		# Export directories setup
		for dir in "${EXPORT_DIR}" "${BACKUP_DIR}"; do
		if [[ ! -d "${dir}" ]]; then
		mkdir -p "${dir}"
		chown nobody:nogroup "${dir}"
		chmod 777 "${dir}"
		echo "LOG: Created directory ${dir}"
		fi
		done
		
		# Environment template verification
		if [[ ! -f ".env" ]]; then
		cp .env.example .env
		echo "WARN: .env configuration initialized!"
		fi
	\end{lstlisting}
	
	\newpage
	
	% ════════════════════════════════════════════════════════════════════════════════
	%  5. НАСТРОЙКА NFS
	% ════════════════════════════════════════════════════════════════════════════════
	\section{Настройка NFS-сервера и клиентов}
	
	Конфигурация NFS-сервера определяется файлом \texttt{/etc/exports}, описывающим
	экспортируемые директории и права доступа для клиентов.
	
	\textbf{Конфиксистема файла \texttt{/etc/exports}:}
	
	\begin{lstlisting}[language=bash, caption={Файл /etc/exports}]
		/exports/shared 172.20.0.0/24(rw,sync,no_subtree_check,fsid=0,root_squash)
		/exports/shared 172.20.0.11(rw,sync,no_root_squash,fsid=0)
	\end{lstlisting}
	
	Строка для подсети \texttt{172.20.0.0/24} включает \texttt{root\_squash} ---
	перемаппинг привилегированных запросов в анонимного пользователя, что снижает риск
	несанкционированного повышения привилегий. Для доверенного клиента
	\texttt{172.20.0.11} (тестовый контейнер) \texttt{no\_root\_squash} разрешает работу
	от root в изолированной среде.
	
	\textbf{Монтирование NFS через Docker Volume.}
	Клиентские контейнеры монтируют NFS-шару напрямую через встроенный драйвер Docker
	Volume без установки дополнительных пакетов в образ:
	
	\begin{lstlisting}[language=yaml, caption={Конфигурация Docker Volume для NFS}]
		volumes:
		nfs_shared:
		driver: local
		driver_opts:
		type: nfs4
		o: "addr=172.20.0.10,rw,sync,nfsvers=4,hard,intr"
		device: ":/exports/shared"
	\end{lstlisting}
	
	\textbf{Контроль жизнеспособности (healthcheck).}
	Для NFS-сервера настроен healthcheck, проверяющий доступность экспорта командой
	\texttt{showmount}:
	
	\begin{lstlisting}[language=yaml, caption={Healthcheck для NFS-сервера}]
		healthcheck:
		test: ["CMD", "showmount", "-e", "localhost"]
		interval: 30s
		timeout: 10s
		retries: 5
		start_period: 20s
	\end{lstlisting}
	
	\newpage
	
	% ════════════════════════════════════════════════════════════════════════════════
	%  6. БЕЗОПАСНОСТЬ
	% ════════════════════════════════════════════════════════════════════════════════
	\section{Системная безопасность и hardening Linux-окружения}
	
	Многоуровневый периметр защиты охватывает сетевой стек ОС хоста, права доступа
	файловой системы и изоляцию контейнеров.
	
	\textbf{Межсетевой экран UFW.}
	Скрипт \texttt{hardening.sh} переводит UFW в режим блокирования входящего трафика по
	умолчанию. Открыты только строго необходимые порты:
	
	\begin{table}[h!]
		\centering
		\small
		\begin{tabularx}{\textwidth}{|l|l|X|}
			\hline
			\textbf{Порт / Протокол} & \textbf{Назначение} & \textbf{Ограничение} \\
			\hline
			22/TCP & SSH --- удалённое администрирование & Только admin-IP \\
			\hline
			2049/TCP & NFS --- внутри storage-net & Подсеть 172.20.0.0/24 \\
			\hline
			3000/TCP & Grafana --- веб-интерфейс & localhost / admin-сеть \\
			\hline
			9090/TCP & Prometheus --- метрики & Только localhost \\
			\hline
		\end{tabularx}
		\caption{Правила межсетевого экрана UFW}
	\end{table}
	
	\textbf{Принцип наименьших привилегий в контейнерах:}
	
	\begin{itemize}
		\item Все контейнеры, кроме \texttt{nfs-server}, запускаются от непривилегированного
		пользователя (\texttt{USER appuser});
		\item \texttt{nfs-server} требует \texttt{privileged: true} для работы с ядерными
		модулями NFS --- это задокументированное исключение;
		\item \texttt{cap\_drop: [ALL]} применяется ко всем клиентским контейнерам;
		\item Корневая ФС клиентов монтируется в режиме \texttt{read\_only: true}.
	\end{itemize}
	
	\newpage
	
	% ════════════════════════════════════════════════════════════════════════════════
	%  7. ТЕСТИРОВАНИЕ, РЕЗЕРВНОЕ КОПИРОВАНИЕ И CI/CD
	% ════════════════════════════════════════════════════════════════════════════════
	\section{Тестирование, резервное копирование и CI/CD}
	
	\textbf{Стратегия интеграционного тестирования.}
	Валидация развёрнутого стенда осуществляется скриптом \texttt{test\_nfs.sh}, который
	проверяет четыре уровня работоспособности:
	
	\begin{enumerate}
		\item \textbf{Сетевая доступность} --- опрашивает экспортируемые директории NFS-сервера
		командой \texttt{showmount -e nfs-server} с таймаутом 60 секунд.
		\item \textbf{Монтирование контейнеров} --- проверяет, что NFS-шара примонтирована в
		каждом клиентском контейнере командой \texttt{df -h | grep nfs}.
		\item \textbf{Совместный доступ} --- записывает файл из \texttt{nfs-client-1} и читает
		его из \texttt{nfs-client-2}, проверяя идентичность содержимого.
		\item \textbf{HA-тест} --- принудительно останавливает \texttt{nfs-server}, ожидает
		перезапуска и проверяет восстановление монтирования.
	\end{enumerate}
	
	\textbf{Резервное копирование.}
	Скрипт \texttt{backup.sh} реализует «горячее» резервное копирование содержимого
	NFS-экспорта без остановки сервиса. Применяется утилита \texttt{rsync} с флагами
	\texttt{--archive --checksum}, обеспечивающими инкрементальное копирование только
	изменённых файлов. Ротация архивов удаляет копии старше 7 дней:
	
	\begin{lstlisting}[language=bash, caption={Ротация архивов в backup.sh}]
		find "${BACKUP_DIR}" -type f -name "*.tar.gz" -mtime +7 -exec rm {} \;
	\end{lstlisting}
	
	\textbf{Конвейер CI/CD (GitHub Actions).}
	Проект интегрирован с облачным конвейером GitHub Actions. Пайплайн
	\texttt{pr-validation.yml} запускается автоматически при открытии Pull Request и
	включает следующие этапы:
	
	\begin{itemize}
		\item \textbf{yamllint} --- проверка синтаксиса всех YAML-файлов проекта;
		\item \textbf{ShellCheck} --- аудит bash-скриптов на логические ошибки и соответствие
		POSIX;
		\item \textbf{Hadolint} --- проверка Dockerfile на оптимизацию слоёв сборки;
		\item \textbf{Динамические тесты} --- развёртывание стенда в эфемерном
		Ubuntu-раннере и запуск \texttt{test\_nfs.sh}.
	\end{itemize}
	
	Слияние кода в ветку \texttt{main} блокируется до успешного прохождения всех этапов
	конвейера.
	
	\newpage
	
	% ════════════════════════════════════════════════════════════════════════════════
	%  ЗАКЛЮЧЕНИЕ
	% ════════════════════════════════════════════════════════════════════════════════
	\section*{Заключение}
	\addcontentsline{toc}{section}{Заключение}
	
	В ходе выполнения данного курсового проекта был успешно спроектирован, развёрнут и
	защищён отказоустойчивый сервис сетевого файлового хранилища на базе протокола NFS в
	контейнеризированной среде под управлением ОС Linux.
	
	Применение методологии Infrastructure as Code и инструментов Docker Compose позволило
	полностью автоматизировать все этапы: подготовку окружения, настройку экспортов,
	монтирование в клиентских контейнерах и проведение интеграционных тестов.
	
	Реализованная многоуровневая система безопасности --- изоляция сети
	\texttt{storage-net}, правила UFW по принципу «запрещено всё», принцип наименьших
	привилегий в контейнерах --- обеспечила высокий уровень защищённости хранимых данных от
	несанкционированного доступа.
	
	Развёрнутый стек мониторинга Prometheus + Grafana с кастомными дашбордами предоставляет
	полную наблюдаемость (observability) работы NFS-сервера в реальном времени. Интеграция
	с CI/CD пайплайном GitHub Actions гарантирует стабильность инфраструктуры на
	протяжении всего жизненного цикла её эксплуатации.
	
	\begin{table}[h!]
		\centering
		\small
		\begin{tabular}{|p{11cm}|c|}
			\hline
			\textbf{Результат} & \textbf{Статус} \\
			\hline
			Развёртывание NFS Replica Set в Docker Compose & $\checkmark$ Выполнено \\
			\hline
			Автоматизация bootstrap.sh и Makefile & $\checkmark$ Выполнено \\
			\hline
			Многоуровневый hardening (UFW, cap\_drop, read\_only) & $\checkmark$ Выполнено \\
			\hline
			Интеграционные тесты и HA-тестирование & $\checkmark$ Выполнено \\
			\hline
			Мониторинг Prometheus + Grafana & $\checkmark$ Выполнено \\
			\hline
			CI/CD пайплайн GitHub Actions & $\checkmark$ Выполнено \\
			\hline
		\end{tabular}
		\caption{Итоговые результаты проекта}
	\end{table}
	
\end{document}